\section{From ranks to LQD coordinates}
\label{sec:model}

\subsection{The speed of a quantile}

Let $Q(u)$ be the quantile function of $X$, with rank $u\in(0,1)$.  If
$U$ is uniform, then $Q(U)$ has the desired law.  Hence probability one is
built into the coordinate $u$, and the only remaining requirement is that
$Q$ increase.

Rank is a universal clock for distributions.  The value $Q(0.01)$ is the
first-percentile return whether the law is symmetric, skewed, or multimodal.
Changing the law changes the returns assigned to the clock, not the clock
itself.  This is the main simplification: normalization of probability is
handled once by the uniform variable rather than reimposed on every candidate
density.

Write $q(u)=Q'(u)$ for the quantile density.  Differentiating
$F_X(Q(u))=u$ gives
\begin{equation}
f_X(Q(u))=\frac1{q(u)}.
\label{eq:reciprocal}
\end{equation}
The function $q$ is the speed with which return space is traversed as rank
increases.  Slow motion packs probability into a narrow return interval;
fast motion thins it.  Since $Q$ is increasing exactly when $q>0$, an
arbitrary real function may be used for $\log q$.

An unbounded distribution forces $q$ to diverge at both endpoints.  Here a
\emph{tail} means the probability lying beyond a remote threshold.  We say
that $X$ has exponential tails, with positive scales $\lambda_-$ and
$\lambda_+$, when
\begin{equation}
F_X(x)\sim C_-e^{x/\lambda_-}\quad(x\to-\infty),
\qquad
1-F_X(x)\sim C_+e^{-x/\lambda_+}\quad(x\to+\infty),
\label{eq:taildefinitions}
\end{equation}
for constants $C_-,C_+>0$.  Their logarithms therefore decay linearly in
$|x|$.  By contrast, a Gaussian tail has log-probability proportional to
$-x^2$: it decays quadratically on the log-probability scale and hence faster
than every exponential tail.  For the exponential class, the leading endpoint
singularities of $q$ are $1/u$ and $1/(1-u)$.  Remove them explicitly:
\begin{equation}
\log q(u)=-\log u-\log(1-u)+g(u).
\label{eq:skeleton}
\end{equation}
The remainder $g$ is bounded for the exponential-tail class considered here.

The subtraction is not arbitrary.  The left relation in
\eqref{eq:taildefinitions} implies
$Q(u)\sim\lambda_-\log u$ as $u\downarrow0$ and
$q(u)\sim\lambda_-/u$.  The right endpoint is analogous.  More precisely, if
$X\sim N(0,s^2)$, then
\[
F_X(x)\sim \frac{s}{|x|\sqrt{2\pi}}e^{-x^2/(2s^2)}
\qquad(x\to-\infty),
\]
with the symmetric formula on the right.  Its quantile density is
$q(u)=s/\phin(\PhiN^{-1}(u))$ and, by the normal tail equivalent,
\[
q(u)\sim\frac{s}{u\sqrt{2\log(1/u)}}\qquad(u\downarrow0).
\]
Its leading $1/u$ singularity is removed, but the remainder still tends to
$-\infty$.  Exponential tails therefore lie inside the LQD chart, whereas
Gaussian tails lie on a different boundary.

Now set
\begin{equation}
z=\logit u=\log\frac{u}{1-u},\qquad
u=\Lambda(z)=\frac1{1+e^{-z}},\qquad
\rho(z)=\Lambda(z)\Lambda(-z).
\label{eq:logodds}
\end{equation}
The function $\Lambda$ is the logistic CDF and $\rho=\Lambda'$ its density.
Because $\dd u/\dd z=u(1-u)$, the singularities in \eqref{eq:skeleton}
cancel:
\begin{equation}
\frac{\dd Q}{\dd z}=q(u)u(1-u)=e^{g(u)}.
\label{eq:speed}
\end{equation}
Thus the entire distribution is generated by integrating a positive,
ordinary function on the log-odds line.

The coordinate also places every percentile at a finite location.  The median
is $z=0$; the tenth and ninetieth percentiles are
$z=\mp\log9$; moving one unit in $z$ multiplies the odds $u/(1-u)$ by $e$.
Meanwhile $\rho(z)\sim e^{-|z|}$, so rank integrals become ordinary integrals
on $\R$ with exponentially decaying weight.  This is why the same coordinate
is convenient both analytically and numerically.

\begin{figure}[htbp]
\centering
\paperfig{\textwidth}{fig_transport.pdf}
\caption{The transport construction for constant speed.  Left: the logistic
score and five percentile marks.  Center: an affine map from log-odds to
log-return, with two strike roots.  Right: the implied-volatility smile priced
from that law.  The first two panels define the distribution; the last merely
expresses it in market coordinates.}
\label{fig:transport}
\end{figure}

\subsection{The LQD slice}

Choose $g$ to be a low-degree polynomial in rank.  The shifted Legendre basis
is convenient because it is orthogonal under the uniform rank measure.  The
constant and linear terms are written as an endpoint chord, leaving higher
modes to describe the body.

\begin{definition}[The LQD slice]
\label{def:lqd}
Fix $N\ge2$ and
$\theta=(L,R,a_2,\ldots,a_N)\in\R^{N+1}$.  Let $Z$ have CDF $\Lambda$ and
density $\rho$.  The LQD slice is the law of
\begin{equation}
\boxed{\;
X=x(Z),\qquad
x(z)=m+\int_0^z e^{g(\Lambda(t))}\,\dd t,\qquad
g(u)=(1-u)L+uR+\sum_{n=2}^N a_nP_n(1-2u),
\;}
\label{eq:lqd}
\end{equation}
where $P_n$ is the Legendre polynomial of degree $n$ and $m$ is chosen so
that $\E[e^X]=1$.
\end{definition}

For every finite coefficient vector, $x'(z)=e^{g(\Lambda(z))}>0$.
Consequently
\begin{equation}
F_X(x(z))=\Lambda(z),\qquad
f_X(x(z))=\frac{\rho(z)}{x'(z)}
=\Lambda(z)\Lambda(-z)e^{-g(\Lambda(z))}>0.
\label{eq:density}
\end{equation}
The quantile is $Q(u)=x(\logit u)$, and \eqref{eq:speed} recovers
\begin{equation}
q(u)=\frac{e^{g(u)}}{u(1-u)}.
\label{eq:lqdform}
\end{equation}
Equations \eqref{eq:lqd} and \eqref{eq:lqdform} are the transport and
log-quantile-density descriptions of the same law.

Since a polynomial $g$ is bounded on $[0,1]$, its exponential is bounded
above and away from zero.  Therefore $x(z)$ tends to $-\infty$ and $+\infty$
with $z$ and is a diffeomorphism of the line.  The construction cannot fold,
stop, or reverse.  Positivity of the density in \eqref{eq:density} is thus a
geometric fact about the coordinate map, not a pointwise inequality checked
after the fit.

The two endpoint values are
\begin{equation}
g(0)=L+\sum_{n\ge2}a_n,\qquad
g(1)=R+\sum_{n\ge2}(-1)^na_n,
\label{eq:endpoints}
\end{equation}
and their exponentials
\begin{equation}
\lambda_-=e^{g(0)},\qquad \lambda_+=e^{g(1)}
\label{eq:speeds}
\end{equation}
are the endpoint speeds.  If
$\bar x(z)=\int_0^ze^{g(\Lambda(t))}\dd t$, then Lipschitz continuity of
$g$ and exponential convergence of $\Lambda$ to its endpoints imply
\begin{equation}
\bar x(z)=b_-+\lambda_-z+O(e^z)\quad(z\to-\infty),
\qquad
\bar x(z)=b_++\lambda_+z+O(e^{-z})\quad(z\to+\infty).
\label{eq:asymptotes}
\end{equation}
The endpoint speeds will therefore govern both tails.

\begin{figure}[htbp]
\centering
\paperfig{\textwidth}{fig_modes.pdf}
\caption{Three Legendre perturbations of the same reference law.  Left:
their action on log-speed $g$.  Center: the response of implied volatility.
Right: the response of the density.  An even mode moves both endpoints in
the same direction; an odd mode moves them in opposite directions.  In all
cases the density remains positive because the transport speed remains an
exponential.}
\label{fig:modes}
\end{figure}

The modes have a direct interpretation.  Raising $g$ stretches returns near
the affected ranks and therefore lowers density there by \eqref{eq:reciprocal}.
Even modes act symmetrically; odd modes transfer shape from one side to the
other; higher modes work on finer rank scales.  The raw coefficients are not
pure body controls, since \eqref{eq:endpoints} couples each of them to one or
both tails.  The endpoint chart introduced in \cref{sec:calibration} removes
that coupling without changing the family.

Legendre polynomials are a coordinate choice, not part of the probability
argument.  Rank integrals use uniform weight on $[0,1]$, under which Legendre
modes are orthogonal.  Their columns remain distinguishable in least-squares
normal equations as order rises.  Monomials span the same space but inherit
Hilbert-matrix conditioning.  Chebyshev modes would also be reasonable; the
endpoint chord would remain useful because it names the two tail controls.

\begin{figure}[htbp]
\centering
\paperfig{0.92\textwidth}{fig_lqd_chart.pdf}
\caption{A fitted SPY December 2026 slice in the model's coordinates.  Left:
the universal boundary skeleton and the fitted $\log q$; their difference is
$g$.  Right: the density produced by the same transport.  The quoted strikes
occupy the central ranks; the continuation to the endpoints is extrapolation.}
\label{fig:lqdchart}
\end{figure}

\subsection{Martingale normalization and the finite-forward wall}

Only a translation remains.  Let
\[
I=\int_{-\infty}^{\infty}e^{\bar x(z)}\rho(z)\,\dd z.
\]

\begin{proposition}[Martingale normalization]
\label{prop:shift}
If $I<\infty$, then $m=-\log I$ is the unique shift for which
$\E[e^X]=1$.  Moreover, $I<\infty$ if and only if $\lambda_+<1$.
\end{proposition}

\begin{proof}
Since $X=m+\bar x(Z)$, $\E[e^X]=e^mI$, giving the unique shift.  At the
right endpoint, \eqref{eq:asymptotes} and $\rho(z)\sim e^{-z}$ make the
integrand asymptotic to a positive constant times
$e^{(\lambda_+-1)z}$; it is integrable exactly when $\lambda_+<1$.  At the
left endpoint it behaves as $e^{(1+\lambda_-)z}$, which is integrable for
every $\lambda_->0$.
\end{proof}

The set of admissible raw parameters is therefore
\begin{equation}
\Theta_N=\{\theta\in\R^{N+1}:\lambda_+<1\}.
\label{eq:admissible}
\end{equation}
It is open and has a single boundary.  A logistic coordinate for
$\lambda_+$ will make that boundary unreachable.

\subsection{The exactly solvable slice}
\label{sec:constspeed}

Set $L=R=\log s$, $a_n=0$, with $0<s<1$.  Then
$x(z)=m+sz$ and $X$ is logistic with scale $s$.  Its moment-generating
function is the beta integral
\begin{equation}
\E[e^{tZ}]
=\int_0^1\left(\frac{u}{1-u}\right)^t\dd u
=\mathrm B(1+t,1-t)
=\Gamma(1+t)\Gamma(1-t)
=\frac{\pi t}{\sin(\pi t)},\qquad |t|<1.
\label{eq:mgf}
\end{equation}
Hence
\[
m=-\log\frac{\pi s}{\sin(\pi s)},\qquad
\operatorname{Var}(X)=\frac{\pi^2s^2}{3}.
\]
Matching a target ATM total variance $w_0$ gives the cold start
\begin{equation}
s=\frac{\sqrt{3w_0}}{\pi}.
\label{eq:coldstart}
\end{equation}
This seed matches variance, not ATM price.  As $s\downarrow0$,
$m=-\pi^2s^2/6+O(s^4)$.  Moreover,
\[
\E[Z^+]=\int_0^\infty z\rho(z)\,\dd z
=\sum_{n\ge1}\frac{(-1)^{n-1}}n=\log2,
\]
where the logistic density is expanded as a geometric series on the positive
half-line.  Thus the LQD ATM call is $s\log2+O(s^2)$.  Matching it to the
small-variance Black price gives
\[
\frac{w_{\rm implied}}{w_0}\longrightarrow
\frac{6(\log2)^2}{\pi}\approx0.918.
\]
The seed is therefore about eight percent low in ATM variance, a known bias
removed by calibration.

\begin{figure}[htbp]
\centering
\paperfig{0.92\textwidth}{fig_exact.pdf}
\caption{The solvable slice as an end-to-end numerical test.  Left: errors in
the computed transport and martingale shift against the closed forms from
\eqref{eq:mgf}; the normalization becomes harder as $s$ approaches the wall.
Right: the variance-matched cold start and its derived small-$s$ bias.}
\label{fig:exact}
\end{figure}
